\documentclass[11pt,a4paper]{article}

% ── Codificación y fuentes ────────────────────────────────────────────────────
\usepackage[utf8]{inputenc}
\usepackage[T1]{fontenc}
\usepackage{lmodern}
\usepackage[spanish]{babel}

% ── Geometría y espaciado ─────────────────────────────────────────────────────
\usepackage[top=2.5cm, bottom=2.5cm, left=2.8cm, right=2.8cm]{geometry}
\usepackage{setspace}
\setstretch{1.15}
\usepackage{parskip}

% ── Tablas ────────────────────────────────────────────────────────────────────
\usepackage{booktabs}
\usepackage{tabularx}
\usepackage{array}
\usepackage{longtable}
\usepackage{enumitem}

% ── Colores y cajas ───────────────────────────────────────────────────────────
\usepackage[dvipsnames]{xcolor}
\usepackage{tcolorbox}
\tcbuselibrary{skins, breakable}

% ── Cabeceras y pies de página ────────────────────────────────────────────────
\usepackage{fancyhdr}
\usepackage{lastpage}
\pagestyle{fancy}
\fancyhf{}
\fancyhead[L]{\small\textbf{Informe de Evaluación} --- Física para Bioanalistas}
\fancyhead[R]{\small Ybianmaris Maribao}
\fancyfoot[C]{\small Página \thepage\ de \pageref{LastPage}}
\renewcommand{\headrulewidth}{0.4pt}

% ── Hiperenlaces ──────────────────────────────────────────────────────────────
\usepackage[colorlinks=true, linkcolor=NavyBlue, urlcolor=NavyBlue]{hyperref}

% ── Paleta de colores personalizados ─────────────────────────────────────────
\definecolor{desempeno_excelente}{HTML}{1A6B2F}
\definecolor{desempeno_bueno}{HTML}{2E5A9C}
\definecolor{desempeno_suficiente}{HTML}{8B6914}
\definecolor{desempeno_deficiente}{HTML}{C0392B}
\definecolor{desempeno_insuficiente}{HTML}{7B241C}
\definecolor{grisFondo}{HTML}{F5F5F5}
\definecolor{azulTitulo}{HTML}{1B2A6B}
\definecolor{verdeFortaleza}{HTML}{1A6B2F}
\definecolor{naranjaMejora}{HTML}{A04000}

% ── Estilos de cajas ─────────────────────────────────────────────────────────
\tcbset{
  cajaTitulo/.style={
    enhanced, breakable,
    colback=azulTitulo!8, colframe=azulTitulo,
    fonttitle=\bfseries\large, coltitle=white,
    attach boxed title to top left={yshift=-2mm, xshift=4mm},
    boxed title style={colback=azulTitulo},
    top=6pt, bottom=6pt, left=8pt, right=8pt,
  },
  cajaFortaleza/.style={
    enhanced, breakable,
    colback=verdeFortaleza!6, colframe=verdeFortaleza!60,
    fonttitle=\bfseries, coltitle=verdeFortaleza,
    top=4pt, bottom=4pt, left=8pt, right=8pt,
  },
  cajaMejora/.style={
    enhanced, breakable,
    colback=naranjaMejora!6, colframe=naranjaMejora!60,
    fonttitle=\bfseries, coltitle=naranjaMejora,
    top=4pt, bottom=4pt, left=8pt, right=8pt,
  },
  cajaRecomendacion/.style={
    enhanced, breakable,
    colback=grisFondo, colframe=gray!50,
    fonttitle=\bfseries, coltitle=black,
    top=4pt, bottom=4pt, left=8pt, right=8pt,
  },
}

% ─────────────────────────────────────────────────────────────────────────────
\begin{document}

% ══ PORTADA ══════════════════════════════════════════════════════════════════
\begin{center}
  {\Large\bfseries\color{azulTitulo} Universidad de Oriente/Núcleo Sucre --- Física para Bioanalistas}\\[4pt]
  {\large Informe de Evaluación Preliminar}\\[12pt]
  \rule{\linewidth}{1.2pt}\\[6pt]
  {\LARGE\bfseries Ybianmaris Maribao}\\[4pt]
  \rule{\linewidth}{0.6pt}
\end{center}

% ══ FICHA TÉCNICA ════════════════════════════════════════════════════════════
\vspace{4pt}
\begin{tcolorbox}[cajaTitulo, title=Datos del informe]
\begin{tabular}{@{}llll@{}}
  \textbf{Estudiante:}  & Ybianmaris Maribao  &
  \textbf{Fecha:}       & 2026-08-08 \\[4pt]
  \textbf{Archivo PDF:} & \texttt{Ybianmaris\_Maribao.pdf}  &
  \textbf{Páginas:}     & 5 \\
\end{tabular}
\end{tcolorbox}

% ══ RESULTADO GLOBAL ═════════════════════════════════════════════════════════
\vspace{6pt}
\begin{center}
  \begin{tcolorbox}[
    enhanced,
    colback=white, colframe=desempeno_bueno,
    width=0.62\linewidth,
    boxrule=2pt,
    halign=center, valign=center,
    top=8pt, bottom=8pt,
  ]
    \centering
    {\large\bfseries Puntaje sugerido}\\[6pt]
    {\Huge\bfseries\color{desempeno_bueno} 7.4/10}\\[6pt]
    {\large Nivel de desempeño: \textbf{\color{desempeno_bueno} Bueno}}
  \end{tcolorbox}
\end{center}

% ══ RESUMEN GENERAL ══════════════════════════════════════════════════════════
\section*{Resumen general}
El estudiante demuestra una sólida comprensión de los principios físicos fundamentales aplicados a bioanálisis en las preguntas que abordó. Sus justificaciones cualitativas son generalmente claras y sus cálculos numéricos son precisos. Sin embargo, se observan áreas de mejora en la profundidad de las explicaciones en algunas preguntas y en la cobertura completa del examen, ya que una pregunta quedó sin responder.

% ══ FORTALEZAS Y ÁREAS DE MEJORA ════════════════════════════════════════════
\vspace{4pt}
\begin{minipage}[t]{0.48\linewidth}
  \begin{tcolorbox}[cajaFortaleza, title={Fortalezas}]
\begin{itemize}[leftmargin=1.5em, itemsep=2pt]
  \item Fuerte comprensión de los principios de termorregulación y calor latente de vaporización (Pregunta 1).
  \item Excelente manejo de la Ley de Laplace y el rol del surfactante pulmonar en la mecánica respiratoria (Pregunta 2).
  \item Correcta aplicación de la Primera Ley de Fick para calcular la eficiencia de difusión en hemodiálisis (Pregunta 3a).
  \item Buen entendimiento de la ósmosis y sus efectos en células sanguíneas ante soluciones hipotónicas (Pregunta 4a).
  \item Precisión en la mayoría de los cálculos numéricos y uso de unidades consistentes en las preguntas resueltas.
  \item Capacidad para justificar respuestas cualitativamente basándose en principios físicos en varias preguntas.
\end{itemize}
  \end{tcolorbox}
\end{minipage}
\hfill
\begin{minipage}[t]{0.48\linewidth}
  \begin{tcolorbox}[cajaMejora, title={Areas de mejora}]
\begin{itemize}[leftmargin=1.5em, itemsep=2pt]
  \item Profundizar en las explicaciones cualitativas y las implicaciones clínicas, especialmente cuando se solicita argumentación detallada (Pregunta 3b y 4b).
  \item Atención a los detalles en las conversiones de unidades finales (Pregunta 3c).
  \item Abordar todas las preguntas del examen para maximizar el puntaje, incluso si la respuesta no es completa (Pregunta 5).
  \item Revisar el tema de capilaridad y fenómenos de superficie.
\end{itemize}
  \end{tcolorbox}
\end{minipage}

% ══ OBSERVACIONES POR PREGUNTA ═══════════════════════════════════════════════
\section*{Observaciones por pregunta}

\begin{longtable}{>{\bfseries}p{2.8cm} p{1.6cm} p{7.0cm} p{2.8cm}}
  \toprule
  \textbf{Pregunta} & \textbf{Puntaje} & \textbf{Evaluación} & \textbf{Errores detectados} \\
  \midrule
  \endhead
  \bottomrule
  \endfoot
    \textbf{Pregunta 1: Termorregulación y Eficiencia de la Sudoración} & 2.0/2.0 & El estudiante responde de manera excelente a todas las partes de esta pregunta. Las justificaciones cualitativas son claras y precisas, explicando correctamente el rol del calor latente de vaporización y el impacto de la humedad relativa. El cálculo numérico es correcto, incluyendo la conversión de unidades. & \textit{---} \\
    \midrule
    \textbf{Pregunta 2: Mecánica Respiratoria y Ley de Laplace} & 2.0/2.0 & La respuesta es impecable. El estudiante demuestra una comprensión profunda de la Ley de Laplace, sus implicaciones en la estabilidad alveolar y el mecanismo de acción del surfactante pulmonar. El cálculo de la diferencia de presión es correcto y bien justificado. & \textit{---} \\
    \midrule
    \textbf{Pregunta 3: Optimización en Hemodiálisis y Primera Ley de Fick} & 1.76/2.0 & En la parte (a), la aplicación de la Primera Ley de Fick y el cálculo del aumento de eficiencia son correctos. En la parte (b), el estudiante identifica el riesgo principal (resistencia mecánica) pero la explicación podría ser más detallada sobre las consecuencias clínicas (ruptura, hemólisis, infección). En la parte (c), el cálculo es correcto en kg/s, pero omite la conversión final a mg/s. & \textit{Explicación incompleta de las consecuencias clínicas del riesgo mecánico (Pregunta 3b).; Omisión de la conversión final de unidades (kg/s a mg/s) en el cálculo (Pregunta 3c).} \\
    \midrule
    \textbf{Pregunta 4: Errores de Infusión Intravenosa y Ósmosis} & 1.64/2.0 & La parte (a) es respondida correctamente, explicando el fenómeno de hemólisis por infusión de agua destilada. Sin embargo, la parte (b) es demasiado concisa; aunque identifica la 'crenación', carece de la explicación detallada de la hipertonicidad, la salida de agua y las implicaciones clínicas (pérdida de flexibilidad, aglomeración, obstrucción). El cálculo de la presión osmótica en la parte (c) es correcto. & \textit{Explicación insuficiente de la crenación y sus implicaciones clínicas (Pregunta 4b).} \\
    \midrule
    \textbf{Pregunta 5: Capilaridad y Toma de Muestras Sanguíneas} & 0.0/2.0 & Esta pregunta quedó completamente en blanco. & \textit{Pregunta sin responder.} \\
    \midrule
\end{longtable}

% ══ ERRORES DETECTADOS ═══════════════════════════════════════════════════════
\section*{Análisis de errores}

\subsection*{Errores conceptuales}
\begin{itemize}[leftmargin=1.5em, itemsep=2pt]
  \item No se detectaron errores conceptuales fundamentales, sino más bien falta de profundidad en algunas explicaciones.
\end{itemize}

\subsection*{Errores procedimentales}
\begin{itemize}[leftmargin=1.5em, itemsep=2pt]
  \item Omisión de la conversión final de unidades (kg/s a mg/s) en el cálculo de la tasa de transporte (Pregunta 3c).
  \item Falta de detalle en la justificación de las consecuencias clínicas en la explicación de la crenación (Pregunta 4b).
\end{itemize}

% ══ RECOMENDACIONES AL ESTUDIANTE ════════════════════════════════════════════
\section*{Retroalimentación para el estudiante}
\begin{tcolorbox}[cajaRecomendacion, title={Dirigido a: Ybianmaris Maribao}]
Tu desempeño en este examen es bueno, demostrando una sólida comprensión de varios principios fundamentales de la física aplicada a bioanálisis, especialmente en termodinámica, mecánica respiratoria y difusión. Tus justificaciones cualitativas en las preguntas 1a, 1b, 2a, 2b y 4a son claras y precisas, y tus cálculos numéricos son mayormente correctos. Para mejorar aún más, te sugiero prestar más atención a la profundidad de tus explicaciones, especialmente cuando se te pide argumentar o justificar. Por ejemplo, en la pregunta 4b, aunque identificaste correctamente la crenación, una explicación más detallada sobre la hipertonicidad y las consecuencias clínicas habría sido beneficiosa. De igual manera, asegúrate de completar todas las conversiones de unidades solicitadas o implícitas en los cálculos. Finalmente, es crucial intentar responder a todas las preguntas del examen. La pregunta 5 quedó en blanco, lo que representa una pérdida significativa de puntos. Incluso si no estás seguro de la respuesta completa, intenta esbozar los principios físicos relevantes o las fórmulas para obtener puntos parciales. Revisa el tema de capilaridad para fortalecer esta área.
\end{tcolorbox}

% ══ NOTA PARA EL PROFESOR ════════════════════════════════════════════════════
\section*{Nota para el profesor}
\begin{tcolorbox}[
  enhanced, breakable,
  colback=yellow!8, colframe=yellow!60!black,
  fonttitle=\bfseries, coltitle=black,
  title={Observaciones para revisión manual},
  top=4pt, bottom=4pt, left=8pt, right=8pt,
]
El estudiante demuestra un buen nivel de comprensión conceptual en las preguntas que abordó (Q1, Q2, Q3, Q4). Las justificaciones cualitativas son generalmente sólidas y los cálculos numéricos son precisos. Los errores son principalmente por falta de detalle en las explicaciones (Q3b, Q4b) o por omisión de una conversión de unidades (Q3c). La pregunta 5 quedó completamente en blanco. El puntaje sugerido de 7.4/10 refleja un desempeño bueno, con áreas claras de mejora en la exhaustividad de las respuestas y la cobertura de todo el examen.
\end{tcolorbox}

% ══ PIE DE INFORME ════════════════════════════════════════════════════════════
\vspace{12pt}
\begin{center}
  \footnotesize\color{gray}
  Informe generado automáticamente por el sistema de evaluación con IA (gemini-2.5-flash) y soporte LaTex. \\
  Este documento es una evaluación \textbf{preliminar} para ser revisado por el profesor antes de ser entregado al 
  estudiante. \\
  Generado el 2026-08-08 \textbullet\ BIOANALISIS Sistema Académico de Evaluación.
  \end{center}

\end{document}
